\section{Introduction}

Open-source projects receive more issues than maintainers can read, and automatic labeling is one of
the most requested forms of help. GitHub now ships an LLM-based labeler as an
Action~\cite{github2025labeler}, and research has moved from fine-tuned classifiers to
LLMs~\cite{colavito2025issue,aracena2025issue}. Benchmarks report how often a labeler is right on a
fixed snapshot~\cite{nlbse24}. A deployed labeler faces different questions. Which labels should it
apply on its own, and which should it leave for a person? Does its promised error rate hold a year
later, or on a repository it has never seen? Does it treat a report differently because the reporter
writes in non-native English or apologizes for being new? And when it explains a decision, is the
explanation what actually drove it?

We study these questions with a triage service built around Qwen2.5-7B, running locally, on the
NLBSE'24 benchmark (3{,}000 issues from five large repositories, 2016--2023) and on 1{,}430 issues the
same repositories received in 2025--2026, after the model's training cutoff. We compare the LLM,
zero-shot and with retrieved examples, against classical classifiers trained on the repositories'
history, and ask:
\begin{description}
\item[RQ1] How do classical and LLM triagers compare in distribution, a year later, and on a new
repository?
\item[RQ2] Can calibrated deferral keep the error among auto-applied labels at a target (10\%), and
does that promise survive drift?
\item[RQ3] Do meaning-preserving changes in how an issue is written (no stack trace, no template, a
newcomer apology, non-native English) change the label?
\item[RQ4] Are the evidence quotes the LLM gives faithful to its decision?
\item[RQ5] Which logged signal notices that the stream of issues has changed?
\end{description}

The answers are mixed in useful ways. A TF-IDF classifier trained on the history is the most accurate
in distribution (76.2\%) but drops to 67.5\% on 2025--26 issues and 60.3\% on a new repository; the
LLM is less accurate in distribution (70.9\%) but stable (75.3\% and 68.0\%). Raw LLM label
probabilities are badly overconfident (fitted temperatures of 8--10). A confidence threshold tuned to
10\% error on calibration data overshoots in every condition (14.9\% in distribution, 20.5\% on a new
repository); a risk-controlled threshold keeps the promise in distribution and a year later but
auto-labels only 7\% of issues. A one-sentence newcomer apology flips 16.2\% of the few-shot LLM's
labels, mostly toward \emph{question}, yet calibrated deferral keeps almost all flips from being
auto-applied. The LLM's evidence quotes are real quotes, but deleting them moves its decision little
more than deleting random text. A simple novelty signal, distance to the closest labeled history
issue, detects a new repository within 11 requests; the model's own confidence rarely notices.

\paragraph{Contributions.} (1) An evaluation of LLM issue triage as a deployed, selective service,
including a post-cutoff 2025--26 set from the benchmark's own repositories; (2) evidence that common
calibration recipes over-promise and that rigorous risk control is affordable only with far more
calibration data; (3) behavioral tests showing a newcomer-apology sensitivity that in-context examples
amplify; (4) a faithfulness test of LLM triage explanations; and (5) monitoring and record/replay
regression tests for the service, all on one laptop.
